% TOP_PROBLEM: {"id": 30007278, "title": "The eleven-twelfths zero-free half-plane", "category_id": 1, "category": {"id": 1, "name": "number_theory", "display_name": "Number Theory", "description": "Properties of integers, prime numbers, Diophantine equations.", "slug": "number-theory", "order_index": 1, "created_at": "2026-07-31T15:26:25.670Z"}, "status": "open"}
% FIRST_POSTED: 2026-10-06
% ATTEMPT_STATUS: solved
% OPENAI_CLAIM: {"schema_version":1,"source_repo":"https://github.com/openai/math","source_commit":"adc7f1241b42e322a6451854ab7e4b4c146bf78a","release_date":"2026-10-06","match":"full","resolution":"proved","independently_reviewed":false,"second_pass":"agreed","families":[{"family":"003","title":"The quasi-Riemann hypothesis","subject":"Number theory","manuscripts":[{"title":"The Quasi-Riemann Hypothesis (alternate 11/12 proof)","pdf_path":"preprints/The-Quasi-Riemann-Hypothesis-October-5-2026/paper2.pdf","readme_path":"preprints/The-Quasi-Riemann-Hypothesis-October-5-2026/README.md","date":"2026-10-05","theorem_ref":"preprints/The-Quasi-Riemann-Hypothesis-October-5-2026/build/paper2.tex :: thm:main"}],"lean":{"doc_path":"lean/docs/003.md","comparators":[],"declarations":[],"covers_main_theorem":false},"reasoning_trace":null}]}
\documentclass[11pt]{article}
\usepackage[margin=25mm]{geometry}
\usepackage{fontspec}
\setmainfont{Latin Modern Roman}
\usepackage{amsmath,amssymb,mathtools}
\usepackage{unicode-math}
\setmathfont{Latin Modern Math}
\usepackage{enumitem}
\usepackage{xurl,hyperref,bookmark}
\hypersetup{colorlinks=true,urlcolor=blue}
\setcounter{secnumdepth}{0}
\setlength{\parindent}{0pt}
\setlength{\parskip}{5pt}
\setlength{\emergencystretch}{3em}
\begin{document}
\section*{The eleven-twelfths zero-free half-plane}

\subsection{OpenAI's claim}
OpenAI presents the following claims relevant to this question, with manuscript numbers following paper order in Sources. In manuscript 1 (family 003; 2026-10-05), the asserted result is that zeta, Dirichlet functions, and finite-order Eisenstein-field Hecke families admit a zero-free region with boundary $11/12$.

\subsection{Outline of the argument}
For manuscript 1 (family 003; 2026-10-05), twisted M\"obius sums are enlarged to a family before Poisson summation. A recursive estimate for the dual mean square, followed by inversion of cube factors, supplies the saving used to exclude a reciprocal-function singularity.

\subsection{Scope relative to this problem}
The result supplies the conclusion in this problem with its displayed hypotheses and quantifiers. The broader statements and applications in the manuscript are kept separate from the particular question above. The supporting result is manuscript 1 (family 003; 2026-10-05).

\subsection{Formal verification}
No Lean formalization is listed at this snapshot for manuscript 1 (family 003; 2026-10-05).

This site has not executed the Lean code. Described statement coverage is not an independent proof check.

\subsection{Status caveat}
The manuscripts were produced by an internal OpenAI model. OpenAI cautions that results without formalization can contain mistakes. This site has not independently checked the proofs; the label records the scope of OpenAI's claim.

\subsection{Sources}
\begin{itemize}
\item OpenAI, \emph{The Quasi-Riemann Hypothesis (alternate 11/12 proof)}, 2026-10-05; citation key \texttt{OAI:The-Quasi-Riemann-Hypothesis-October-5-2026}. \href{https://github.com/openai/math/blob/adc7f1241b42e322a6451854ab7e4b4c146bf78a/preprints/The-Quasi-Riemann-Hypothesis-October-5-2026/paper2.pdf}{PDF (GitHub)}; \href{https://github.com/openai/math/raw/adc7f1241b42e322a6451854ab7e4b4c146bf78a/preprints/The-Quasi-Riemann-Hypothesis-October-5-2026/paper2.pdf}{PDF (direct)}; \href{https://github.com/openai/math/blob/adc7f1241b42e322a6451854ab7e4b4c146bf78a/preprints/The-Quasi-Riemann-Hypothesis-October-5-2026/README.md}{Manuscript page and citation}; \href{https://github.com/openai/math/blob/main/preprints/The-Quasi-Riemann-Hypothesis-October-5-2026/paper2.pdf}{Latest version}.
\item Family 003: \href{https://github.com/openai/math/blob/adc7f1241b42e322a6451854ab7e4b4c146bf78a/lean/docs/003.md}{Lean scope}.
\item \href{https://github.com/openai/math/blob/adc7f1241b42e322a6451854ab7e4b4c146bf78a/CONTENTS.md}{OpenAI catalogue entry}.
\end{itemize}
\end{document}
